\paragraph{SM 签名的代数刚性、相邻层禁激活与进位吸收}\label{par:bdry-tower-sm-rigidity-adjacent-carry}
本段将边界三层塔的 \(D=12\) 计数接口（\S\ref{subsec:bdry-tower-zeck-gut}）与 Zeckendorf--Fibonacci 签名字典（表 \ref{tab:zeckendorf_gut_signatures}）进一步提升为三条可引用的结构性断言：一条关于 \(D=12\) 的紧李代数刚性，一条关于相邻层禁激活导致的 strict 包含不可能性，一条关于 \(m=9\) 激发在审计签名层面的进位吸收。

\begin{proposition}[Zeckendorf 原子性下的 SM 规范代数刚性]\label{prop:sm-zeckendorf-lie-algebra-rigidity}
在“规范玻色子总数 \(D\) 由边界塔层激活给出”且“无冗余激活”由 Zeckendorf 唯一分解审计的口径下，
若 \(D=12\)，则其 Zeckendorf 分解唯一为
\[
12=F_{6}+F_{4}+F_{2},
\qquad
\Sigma(\mathfrak g_{\mathrm{SM}})=\{8,6,4\}
\]
（表 \ref{tab:zeckendorf_gut_signatures}）。
进一步地，若要求“各因子维数与 Zeckendorf 项逐项对齐”（即：每个被激活层对应一个不可再分的紧简单/阿贝尔因子维数槽位），
则任一与该审计兼容的紧李代数 \(\mathfrak g\) 必同构于
\[
\boxed{\ \mathfrak g\ \cong\ \mathfrak u(1)\ \oplus\ \mathfrak{su}(2)\ \oplus\ \mathfrak{su}(3)\ }.
\]
并且该分解在“逐项对齐”的意义下唯一。
\leanverified{sm\_zeckendorf\_twelve}
\leanverified{sm\_boundary\_count}
\leanverified{sm\_zeckendorf\_no\_adjacent}
\leanverified{paper\_sm\_zeckendorf\_lie\_algebra\_rigidity}
\end{proposition}
\begin{proof}
由 Zeckendorf 唯一性，\(12\) 的分解只能为 \(F_6+F_4+F_2=8+3+1\)，且禁止重复项与相邻项。
在紧李代数分类中，维数为 \(1\) 的紧李代数必为 \(\mathfrak u(1)\)；
维数为 \(3\) 的紧简单李代数在同构意义下唯一（\(\mathfrak{su}(2)\cong\mathfrak{so}(3)\)）；
维数为 \(8\) 的紧简单李代数亦唯一（\(\mathfrak{su}(3)\)）。
在“逐项对齐”的口径下三项维数槽位只能分别由上述三因子承担，从而得到所示直和分解，并且唯一性随之成立。证毕。
\end{proof}

% AUTO-GENERATED by scripts/exp_foldbin4_groupoid_aut_pi1_torsion_exponent.py
\begin{equation}\label{eq:foldbin4_groupoid_aut_pi1_torsion_exponent}
\begin{aligned}
\#\{d^{\mathrm{bin}}_4=d\}&:\quad 1:3,\ 2:2,\ 3:3,\\
\exp\bigl(\pi_1(\mathrm{Aut}(\CC[\mathcal{R}_4^{\mathrm{bin}}])^0)\bigr)&=\operatorname{lcm}\{d^{\mathrm{bin}}_4(w):\ w\in X_4\}=6.
\end{aligned}
\end{equation}


\begin{theorem}[SM 三层签名上的共同循环识别：阶数必整除 \texorpdfstring{$6$}{6}]\label{thm:sm-signature-common-cyclic-identification-divides-6}
令 \(\mathcal G_m:=\mathrm{Aut}(\CC[\mathcal{R}_m^{\mathrm{bin}}])^0\) 为 dyadic fold 的二进制关系群胚代数之连通 \(\ast\)-自同构群（附录 \ref{subsec:foldbin-groupoid-aut-lie-selection}），并记其基本群挠指数为
\(
\exp(\pi_1(\mathcal G_m))
=\operatorname{lcm}\{d_m^{\mathrm{bin}}(w):w\in X_m\}
\)
（命题 \ref{prop:foldbin-groupoid-aut0-pi1-torsion-exponent}）。
若存在循环群 \(\ZZ/n\ZZ\) 使对每个 \(m\in\Sigma(\mathfrak g_{\mathrm{SM}})=\{4,6,8\}\) 都存在群单射
\[
\ZZ/n\ZZ\hookrightarrow \pi_1(\mathcal G_m),
\]
则必有
\[
n\mid 6.
\]
并且若进一步要求 \(2\mid n\) 且 \(3\mid n\)，则唯一可能为 \(n=6\)。
\leanverified{paper\_sm\_signature\_common\_cyclic}
\leanverified{paper\_sm\_signature\_common\_cyclic\_identification\_divides\_6}
\end{theorem}
\begin{proof}
由推论 \ref{cor:foldbin-groupoid-aut0-cyclic-injection-criterion}，\(\ZZ/n\ZZ\hookrightarrow \pi_1(\mathcal G_m)\) 当且仅当 \(n\mid \exp(\pi_1(\mathcal G_m))\)。
将其应用于 \(m=4\) 并结合 \eqref{eq:foldbin4_groupoid_aut_pi1_torsion_exponent} 中 \(\exp(\pi_1(\mathcal G_4))=6\)，得到 \(n\mid 6\)。
若再要求 \(2\mid n\) 且 \(3\mid n\)，则由 \(n\mid 6\) 立得 \(n=6\)。证毕。
\end{proof}

\begin{conjecture}[挠谱约束下的 SM 全局形态抽取]\label{conj:sm-global-form-extraction-from-torsion-spectrum}
令 \(\mathcal G_m=\mathrm{Aut}(\CC[\mathcal{R}_m^{\mathrm{bin}}])^0\)。
存在一类自然的闭正规子群 \(N_m\triangleleft \mathcal G_m\)（仅依赖于 dyadic fold 的可审计数据），使得当 \(m\in\Sigma(\mathfrak g_{\mathrm{SM}})=\{4,6,8\}\) 时，
商群
\[
\mathcal G_m^{\mathrm{obs}}:=\mathcal G_m/N_m
\]
在同构意义下稳定为紧连通 Lie 群
\[
\frac{SU(3)\times SU(2)\times U(1)}{\ZZ_6}.
\]
其中 \(\ZZ_6\) 由定理 \ref{thm:sm-signature-common-cyclic-identification-divides-6} 的共同循环识别阶数与 window-\(6\) 的 \(\ZZ_2/\ZZ_3\) 源拼接（推论 \ref{cor:bdry-global-z6-quotient}）共同锁定。
\end{conjecture}

\begin{theorem}[window-\(6\) 的平方余量刚性：\texorpdfstring{$F_{m+2}-12=F_{m-2}^2$}{F_{m+2}-12=F_{m-2}^2} 的唯一解]\label{thm:sm-square-residual-rigidity-m6}
在边界塔计数口径下（\(|X_m|=F_{m+2}\)、\(|X_m^{\mathrm{bdry}}|=F_{m-2}\)，见 \S\ref{subsec:bdry-tower-zeck-gut}），考虑丢番图方程
\[
F_{m+2}-12=\bigl|X_m^{\mathrm{bdry}}\bigr|^2=F_{m-2}^2,
\qquad m\ge 6.
\]
则其在整数 \(m\ge 6\) 上唯一解为 \(m=6\)；并且在该解处成立严格恒等式
\[
|X_6|=F_8=21=12+3^2
=12+\bigl|X_6^{\mathrm{bdry}}\bigr|^2,
\qquad |X_6^{\mathrm{bdry}}|=F_4=3.
\]
\leanverified{sm\_square\_residual\_rigidity\_m6}
\leanverified{fib\_sq\_gt\_fib\_shift}
\leanverified{paper\_sm\_square\_residual\_rigidity\_m6}
\end{theorem}
\begin{proof}
当 \(m=6\) 时，\(|X_6|=F_8=21\) 且 \(|X_6^{\mathrm{bdry}}|=F_4=3\)，故 \(21-12=9=3^2\)，从而 \(m=6\) 为解。
当 \(m=7\) 时，\(F_9-12=34-12=22\neq 25=F_5^2\)，故 \(m=7\) 非解。
\par
对任意 \(m\ge 8\)，令 \(n:=m-2\ge 6\)。
由 Fibonacci 递推可得
\[
F_{n+4}=3F_{n+1}+2F_n<3(2F_n)+2F_n=8F_n,
\]
其中不等式使用 \(F_{n+1}<2F_n\)（\(n\ge 2\)）。
又因 \(n\ge 6\Rightarrow F_n\ge 8\)，故 \(F_n^2\ge 8F_n>F_{n+4}\)，即
\[
F_{m-2}^2=F_n^2>F_{n+4}=F_{m+2}\ge F_{m+2}-12,
\]
从而方程不可能成立。综上唯一解为 \(m=6\)。证毕。
\end{proof}

\begin{corollary}[相邻层禁激活与 strict 包含不可能性]\label{cor:zeckendorf-no-adjacent-strict-inclusion-impossible}
对任意候选紧李代数 \(\mathfrak g\)，其分辨率签名 \(\Sigma(\mathfrak g)=\{k+2:k\in S\}\) 由 Zeckendorf 指标集合 \(S\) 给出且 \(S\) 不含相邻指标（\S\ref{subsec:bdry-tower-zeck-gut}）。
因此对任意 \(m\) 成立
\[
m\in\Sigma(\mathfrak g)\ \Longrightarrow\ m\pm 1\notin\Sigma(\mathfrak g).
\]
特别地，若取 strict 口径要求
\(
\Sigma(\mathfrak g)\supseteq \Sigma(\mathfrak g_{\mathrm{SM}})=\{4,6,8\}
\)
（表 \ref{tab:zeckendorf_gut_signatures}），则必有
\(
\{5,7\}\cap \Sigma(\mathfrak g)=\varnothing
\)。
从而任何签名包含 \(m=5\) 或 \(m=7\) 的候选不可能在签名层面 strict 包含 SM；例如 \(E_6\) 的签名为 \(\Sigma(E_6)=\{12,10,5\}\)（表 \ref{tab:zeckendorf_gut_signatures}），故不满足 strict 包含。
\leanverified{zeckendorf\_no\_adjacent\_strict\_inclusion}
\leanverified{paper\_zeckendorf\_no\_adjacent\_strict\_inclusion\_impossible}
\end{corollary}
\begin{proof}
由 Zeckendorf 条件“指标集合不含相邻指标”，若 \(m=k+2\in\Sigma(\mathfrak g)\)，则 \(k\pm 1\notin S\)，等价于 \(m\pm 1\notin \Sigma(\mathfrak g)\)。
将其应用于 \(\Sigma(\mathfrak g)\supseteq\{4,6,8\}\) 立得 \(\{5,7\}\cap\Sigma(\mathfrak g)=\varnothing\)；最后一条由表 \ref{tab:zeckendorf_gut_signatures} 直接核验。证毕。
\end{proof}

\begin{proposition}[相邻层禁激活下的 uplift 残余中介：\texorpdfstring{$m=7$}{m=7} 的二阶差分回注]\label{prop:sm-uplift-residual-mediator-m7}
\leanverified{paper\_sm\_uplift\_residual\_mediator\_m7}
在 SM 背景签名 \(\Sigma(\mathfrak g_{\mathrm{SM}})=\{4,6,8\}\) 已固定且相邻层禁激活成立的前提下，
尽管 \(m=7\notin\Sigma(\mathfrak g_{\mathrm{SM}})\)，boundary uplift 的三层差值数据仍在算术上内生编码出 window-\(6\) 的 \(\delta=34\) 二层差值：
由表 \ref{tab:foldbin_boundary_lift_m6_m8} 的 \(m=7\) 行 \(\{0,55,89\}\) 可得二阶差分残余
\[
89-55=34,
\]
从而在跨分辨率审计口径下，\(m=7\) 作为被签名排除的相邻层仍提供一条不可由符号层直接消去的“残余差分通道”（定理 \ref{thm:bdry-uplift-second-difference-residual-law}）。
\end{proposition}
\begin{proof}
该恒等式为表 \ref{tab:foldbin_boundary_lift_m6_m8} 所列 uplift-delta 的直接算术计算，
并由定理 \ref{thm:bdry-uplift-second-difference-residual-law} 统一表述为二阶差分残余律。证毕。
\end{proof}

\begin{corollary}[协议激发层与审计签名层的分离：\texorpdfstring{$m=9$}{m=9} 的进位吸收]\label{cor:zeckendorf-carry-absorption-m9}
在 SM 背景签名 \(\Sigma(\mathfrak g_{\mathrm{SM}})=\{4,6,8\}\) 已固定的前提下，
任何试图额外激活 \(m=9\)（对应新增 \(F_7=13\) 个边界模态）的构造在 Zeckendorf 审计签名层面均不可保持为”非冗余新层”：
一旦把层集合正规化为 Zeckendorf 规范形，\(m=8\) 与 \(m=9\) 将触发唯一的进位吸收
\[
F_{7}+F_{6}=F_{8},
\qquad\text{等价地}\qquad
\{8,9\}\ \leadsto\ \{10\}.
\]
因此，\(m=9\) 只能作为协议层的瞬时激发存在；在 \(\Sigma(\cdot)\) 的最小非冗余审计意义下，最小上跳必然落在 \(m=10\)。
\leanverified{zeckendorf\_carry\_absorption\_m9}
\leanverified{paper\_zeckendorf\_carry\_absorption\_m9}
\end{corollary}
\begin{proof}
\(\Sigma(\mathfrak g)\) 的定义要求 Zeckendorf 指标集合不含相邻指标（\S\ref{subsec:bdry-tower-zeck-gut}）。
由于 SM 背景包含 \(m=8\)（即包含 \(F_6\) 项），若再加入 \(m=9\)（即加入 \(F_7\) 项），则出现相邻指标 \(6,7\)，违反规范形条件。
Zeckendorf 正规化在该相邻对上唯一的消元为 \(F_6+F_7=F_8\)，对应层级替换 \(\{8,9\}\leadsto\{10\}\)。
故在审计签名层面 \(m=9\) 不形成独立新层，结论成立。证毕。
\end{proof}

\begin{corollary}[共振梯—四步平移—Zeckendorf 约束下的极小三元组选择律]\label{cor:sm-minimal-triple-selection-law}
\leanverified{paper\_sm\_minimal\_triple\_selection\_law}
在边界词塔口径下，若要求“至少包含两种非阿贝尔共振层”在分辨率签名上表现为
\(
\{6,8\}\subseteq \Sigma(\mathfrak g)
\)
（参见推论 \ref{cor:fib-lie-resonance-scarcity-su2-su3} 与注 \ref{rem:fib-lie-resonance-shift-by-4} 的 \(m\mapsto m+4\) 平移对应），并且层集合按 Zeckendorf 规范形取“无相邻指标、无冗余激活”，
则 \(\Sigma(\mathfrak g)\) 的算术极小解唯一为
\[
\Sigma(\mathfrak g)=\Sigma(\mathfrak g_{\mathrm{SM}})=\{4,6,8\},
\qquad
D=\sum_{m\in\Sigma(\mathfrak g)}|X_m^{\mathrm{bdry}}|=1+3+8=12.
\]
从而在逐项对齐口径下，\(\mathfrak g\) 的同构型被强制为
\(
\mathfrak u(1)\oplus\mathfrak{su}(2)\oplus\mathfrak{su}(3)
\)
（命题 \ref{prop:sm-zeckendorf-lie-algebra-rigidity}）。
\end{corollary}
\begin{theorem}[SM Zeckendorf 槽位的内涨落筛选与坏骨架三分]\label{distill:connes_spectral_triples_and_chamseddine_connes_spectral_action-inner_fluctuation_gauge_family-sm-zeckendorf-inner-fluctuation-bad-skeleton-trichotomy}
令
\[
\Sigma_{\mathrm{SM}}:=\{4,6,8\}=\Sigma(\mathfrak g_{\mathrm{SM}})
\]
为命题 \ref{prop:sm-zeckendorf-lie-algebra-rigidity} 中的 SM Zeckendorf 签名槽位，并分别记
\[
d_4=1,\qquad d_6=3,\qquad d_8=8.
\]
对每个 \(m\in\Sigma_{\mathrm{SM}}\)，设
\(\mathcal A_m:=\CC[\mathcal R_m^{\mathrm{bin}}]\)，并给定忠实有限维 \(\ast\)-表示
\(\pi_m:\mathcal A_m\to\End(H_m)\) 及自伴更新算子 \(D_m\)。以内涨落规范抽取算子作为有限接口，定义实自伴一形式空间
\[
\Omega^1_{m,{\rm sa}}(D_m):=\Span_{\mathbb R}\{T+T^*:T=\pi_m(a)[D_m,\pi_m(b)],\ a,b\in\mathcal A_m\}.
\]
设
\[
\mathcal Z_m:=\Omega^1_{m,{\rm sa}}(D_m)\cap \pi_m(\mathcal A_m)',
\qquad
Q_m:\Omega^1_{m,{\rm sa}}(D_m)\longrightarrow
\overline\Omega^1_m:=\Omega^1_{m,{\rm sa}}(D_m)/\mathcal Z_m
\]
为去除可观测平凡一形式后的商映射。又假设 \(\mathcal G_m=\Aut(\CC[\mathcal R_m^{\mathrm{bin}}])^0\) 在 \(\overline\Omega^1_m\) 上有良定义的线性输运 \(\overline\Gamma_{m,g}\)，且对每个 \(u\in U(\mathcal A_m)^0\) 的内自同构有
\[
\overline\Gamma_{m,\Ad u}(Q_m(A))=Q_m\bigl(\pi_m(u)A\pi_m(u)^*\bigr).
\]
设一个槽位对齐候选由紧李代数分解
\[
\mathfrak h=\bigoplus_{m\in\Sigma_{\mathrm{SM}}}\mathfrak h_m,\qquad
\dim\mathfrak h_m=d_m,
\]
以及实线性映射
\(\Phi_m:\mathfrak h_m\to\Omega^1_{m,{\rm sa}}(D_m)\) 给出。记
\[
W_m:=Q_m\Phi_m(\mathfrak h_m),
\]
并定义三个 obstruction 空间
\[
K_m:=\ker(Q_m\Phi_m),
\]
\[
E_m:=\left(\sum_{g\in\mathcal G_m}\overline\Gamma_{m,g}(W_m)\right)/W_m,
\]
其中分子表示由所有输运像生成的有限维实线性子空间，以及
\[
P_m:=\bigl(W_m+\Theta_m\bigr)/W_m,\qquad
\Theta_m:=\Span_{\mathbb R}\{Q_m(\pi_m(u)[D_m,\pi_m(u^*)]):u\in U(\mathcal A_m)^0\}.
\]
则下列条件等价：
\begin{enumerate}
  \item 对所有 \(m\in\Sigma_{\mathrm{SM}}\)，有
  \[
  K_m=0,\qquad E_m=0,\qquad P_m=0.
  \]
  \item 每个 \(W_m\) 是一个 \(d_m\)-维可观测内涨落槽位，且同时满足：\(W_m\) 对 \(\mathcal G_m\) 的自同构输运稳定；并且对所有 \(u\in U(\mathcal A_m)^0\)，仿射规范变换
  \[
  \mathcal T_{m,u}(A):=\pi_m(u)A\pi_m(u)^*+\pi_m(u)[D_m,\pi_m(u^*)]
  \]
  在商空间中把 \(W_m\) 保持到自身。
\end{enumerate}
因此，任一满足 Zeckendorf 维数计数 \(1+3+8=12\) 的 group\_unification 候选，若不能通过内涨落与自同构筛选，则必且只必出现以下三类坏骨架之一：存在某个 \(m\) 使 \(K_m\ne0\) 的可观测一形式坍缩；或存在某个 \(m\) 使 \(E_m\ne0\) 的自同构输运逃逸；或存在某个 \(m\) 使 \(P_m\ne0\) 的纯规范内涨落外逃。
\end{theorem}
\begin{proof}
先注意，若 \(u\in U(\mathcal A_m)^0\)，则
\[
\pi_m(u)[D_m,\pi_m(u^*)]
=\pi_m(u)D_m\pi_m(u)^*-D_m
\]
为自伴算子，且属于 \(\Omega^1_{m,{\rm sa}}(D_m)\)。故 \(\Theta_m\) 定义良好。

设首先 \(K_m=E_m=P_m=0\)。由 \(K_m=0\) 与 \(\dim\mathfrak h_m=d_m\)，映射 \(Q_m\Phi_m\) 将 \(\mathfrak h_m\) 同构到 \(W_m\)，故 \(\dim W_m=d_m\)。由 \(E_m=0\)，所有 \(\overline\Gamma_{m,g}(W_m)\) 均包含于 \(W_m\)；取 \(g^{-1}\) 再用同样结论，得 \(\overline\Gamma_{m,g}(W_m)=W_m\)。因此 \(W_m\) 对 \(\mathcal G_m\) 的自同构输运稳定。由 \(P_m=0\)，有 \(\Theta_m\subseteq W_m\)。于是对任意 \(A\in\Omega^1_{m,{\rm sa}}(D_m)\) 且 \(Q_m(A)\in W_m\)，利用内自同构的输运兼容性，得到
\[
Q_m(\mathcal T_{m,u}(A))
=\overline\Gamma_{m,\Ad u}(Q_m(A))+Q_m(\pi_m(u)[D_m,\pi_m(u^*)])\in W_m.
\]
故仿射规范变换在商空间中保持 \(W_m\)。

反过来，若每个 \(W_m\) 是 \(d_m\)-维可观测槽位，则 \(Q_m\Phi_m\) 从 \(d_m\)-维空间到 \(d_m\)-维像空间无核，故 \(K_m=0\)。若 \(W_m\) 对所有 \(\mathcal G_m\)-输运稳定，则
\[
\sum_{g\in\mathcal G_m}\overline\Gamma_{m,g}(W_m)=W_m,
\]
从而 \(E_m=0\)。最后，由仿射规范稳定性应用于 \(A=0\)，得
\[
Q_m(\pi_m(u)[D_m,\pi_m(u^*)])=Q_m(\mathcal T_{m,u}(0))\in W_m
\]
对所有 \(u\in U(\mathcal A_m)^0\) 成立，故 \(\Theta_m\subseteq W_m\)，即 \(P_m=0\)。

还需验证上述仿射变换确为内部规范作用，而非逐点偶然闭合。对 \(u,v\in U(\mathcal A_m)^0\)，直接计算
\[
\mathcal T_{m,u}(\mathcal T_{m,v}(A))
=\pi_m(uv)A\pi_m(uv)^*+\pi_m(uv)[D_m,\pi_m((uv)^*)]
=\mathcal T_{m,uv}(A),
\]
其中使用
\[
[D_m,\pi_m(v^*u^*)]=[D_m,\pi_m(v^*)]\pi_m(u^*)+\pi_m(v^*)[D_m,\pi_m(u^*)].
\]
因此 \(P_m=0\) 与 \(E_m=0\) 排除的是整个连通酉规范轨道的逃逸，而不只是某个生成元的线性误差。

最后，命题 \ref{prop:sm-zeckendorf-lie-algebra-rigidity} 已将 \(D=12\) 的槽位维数刚性固定为 \(1,3,8\)。由已证等价性，一个槽位对齐候选未能通过内涨落规范抽取与自同构筛选，当且仅当至少一个槽位的上述三个 obstruction 空间之一非零。这三类分别对应可观测一形式维数坍缩、自同构输运不稳定、以及纯规范内涨落不在命名槽位中闭合。证毕。
\end{proof}

\begin{proof}
由推论 \ref{cor:fib-lie-resonance-scarcity-su2-su3}，稳定类型共振梯在最小窗口上对 \(\dim\mathfrak{su}(2)=3\) 与 \(\dim\mathfrak{su}(3)=8\) 给出最早的非阿贝尔共振点；
由注 \ref{rem:fib-lie-resonance-shift-by-4}，在边界塔口径下该两点对应边界层 \(m=6,8\)，故“至少两种非阿贝尔共振层”迫使 \(\{6,8\}\subseteq\Sigma(\mathfrak g)\)。
Zeckendorf 规范形禁止相邻层同时出现，因而 \(m=7,9\) 皆不可与 \(\{6,8\}\) 共存。
在允许加入的其余层中，最小者为 \(m=4\)（贡献 \(|X_4^{\mathrm{bdry}}|=F_2=1\)），从而得到唯一的极小签名 \(\{4,6,8\}\)；
由 \(|X_m^{\mathrm{bdry}}|=F_{m-2}\) 得 \(D=1+3+8=12\)。
最后由命题 \ref{prop:sm-zeckendorf-lie-algebra-rigidity} 得到 \(\mathfrak g\) 的同构型与唯一性。证毕。
\end{proof}
